\documentclass[11pt, a4paper]{article}

\usepackage{ctex}
\usepackage{geometry}
\usepackage{titlesec}
\usepackage{enumitem}
\usepackage{setspace}
\usepackage{hyperref}


\title{Script}
\author{黄浩宇}
\date{\today}

\begin{document}
	
	\maketitle
	\newpage
	\tableofcontents
	
	\newpage
	
	\section{【人物表】}
	
	\begin{itemize}[leftmargin=2em, itemsep=0.2em]
		\item \textbf{校长}：学校最高决策者，极力推行形式化改革
		\item \textbf{开明老师}：对形式化教育持有异议的老师
		\item \textbf{教导主任（待出场）}：学校管理层成员
		\item \textbf{教育部部长}：上级视察领导，注重形式与政绩
		\item \textbf{秘书}：部长的随行人员，负责处理突发情况
		\item \textbf{学生A、B、C}：面对视察表现各异的学生
		\item \textbf{群众领导、老师}：随行听会、顺从形式的人员
		\item \textbf{群众学生}：机械化配合学校安排的学生群体
	\end{itemize}
	
	\newpage
	
	\section{【剧本正文】}
	
		\subsection{开场}
			\textbf{△ 画面/音效：} 屏幕上出现一只巨大的鲨鱼充满血腥的嘴巴和尖牙，并且播放凄厉的惨叫声。
			
			\vspace{1em}
		\ \ \
		\subsection{第1幕}
	
			\subsubsection{场景信息}
				\textbf{场景：} 会议室 \quad \textbf{时间：} 白天 \quad \textbf{人物：} 校长、开明老师、群众领导、群众老师
			
			\subsubsection{剧情内容}
				\textbf{△ 场景布置：} 校长站在讲台上给老师和领导们开会。
				
				\vspace{0.5em}
				
				\textbf{校长：} 最近教育部要求提高学生综合素养，不要只会死读书，大家赞不赞同？
				
				\textbf{△ 动作/音效：} 全场所有人奏响机器一般的掌声，非常整齐、利落。掌声响完三声以后，立刻停止，整齐划一。
				
				\textbf{校长：} 为了全面提高学生的综合素养，我们将星期六由放假改为上课，开设各种课外兴趣培训班，好不好？
				
				\textbf{△ 动作/音效：} 全场所有人奏响机器一般的掌声，非常整齐、利落。掌声响完三声以后，立刻停止，整齐划一。
				
				\textbf{开明老师：} （站起来） 这样会占用学生的课余时间。
				
				\textbf{校长：} 为了抢救学生而违背教育理念，好不好笑？
				
				\textbf{△ 动作/音效：} 全场礼歌响起，非常整齐划一的笑声：“哈哈哈”，并且整齐划一地结束。
				
				\textbf{校长：} 散会。
				
				\textbf{△ 动作：} 所有老师迈着机器人一样的步子走出会场，全身几乎不动，只有手和腿僵硬地摆动，一步一步走出教室。
				
				\vspace{1em}
		\ \ \
		\subsection{第2幕}
	
			\subsubsection{场景信息}
				\textbf{场景：} 教室 \quad \textbf{时间：} 星期六，白天 \quad \textbf{人物：} 老师、群众学生
	
			\subsubsection{剧情内容}
				\textbf{△ 场景布置：} 教室顶上的电子钟显示：星期六。黑板上写着：电路设计兴趣小组。所有学生像机器人一样坐在桌子前，桌子上摆放着一块电路板，讲台上放着一个电子报时器。
				
				\vspace{0.5em}
				
				\textbf{△ 动作/音效：} 突然报时器“滴”的一声响起，所有学生整齐划一地把电路通电。电路开始工作，亮起一盏小黄灯。
				
				\textbf{老师：} 大家有没有感觉到自己的素质提高了？
				
				\textbf{所有学生：} （异口同声） 对。
				
				\textbf{△ 动作/音效：} 突然报时器“滴”的一声响起，所有学生马上站起来，迈着机器人一样的步伐，和老师一起离开教室。
				
				\textbf{△ 动作：} 黑板上的字自动改名为：兴趣小组结束。
				
				\textbf{△ 画面/音效：} 相机快门声音响起。
				
				\textbf{△ 画面显示/新闻内容：} 学校网站上登报一则新闻：“XX高中在周末举办兴趣小组，学生的综合素质普遍提高，符合国家先进教育理念，使得我们高中成为全市领先的独一档的高中！”
				
				\vspace{1em}
		\ \ \
		\subsection{第3幕}
	
			\subsubsection{场景信息}
				\textbf{场景：} 学校走廊 / 校园内 \quad \textbf{时间：} 白天 \quad \textbf{人物：} 教育部部长、校长、秘书、学生A、学生B、学生C
			
			\subsubsection{剧情内容}
				\textbf{△ 场景布置：} 教育部部长在校长等人的陪同下，亲自到学校视察。
				
				\vspace{0.5em}
				
				\textbf{△ 动作：} 部长迎面拦住走过来的学生A，进行现场采访。
				
				\textbf{教育部部长：} 你觉得你们学校的周六兴趣班怎么样？
				
				\textbf{学生A：} 不怎么样，很无聊，听不懂，也不想做。
				
				\textbf{△ 动作：} 部长的脸色马上转为沉重。他向旁边的秘书招了招手，秘书马上过来把学生A拉走。
				
				\textbf{△ 动作：} 部长马上把脸色改回微笑。
				
				\textbf{△ 动作：} 这时又走过来学生B。
				
				\textbf{校长：} （上前问学生B） 你觉得你们学校的周六兴趣班怎么样？
				
				\textbf{学生B：} 不怎么样，很无聊，听不懂，也不想做。
				
				\textbf{△ 动作：} 部长的脸色马上再次转为沉重，向旁边的秘书招招手，秘书马上把学生B拉走。
				
				\textbf{△ 动作：} 接着又走过来学生C。
				
				\textbf{校长：} （问学生C） 你觉得你们学校的周六兴趣班怎么样？
				
				\textbf{学生C：} 这个活动非常有意义，不仅提高了我们的综合素质水平，而且开拓了视野，真是令学生受益匪浅。
				
				\textbf{△ 动作：} 部长马上开怀大笑，拍了拍学生C的肩膀。
				
				\textbf{教育部部长：} 加把劲，继续努力！
				
				\textbf{△ 画面/新闻内容：} 随后，本市教育部发布新闻：“本市教育水平已经领先全省，全市学生综合素质普遍提高，学生对教育部的举措认可度极高，已成为全省的教育先进市！”
				
\end{document}